% !TeX root = manuscript.tex
Text analysis in economics has traditionally proceeded by mapping documents into structured numerical representations, such as dictionaries that count curated terms and phrases, bag-of-words or TF--IDF features that count word occurrences, and topic or factor models that compress word counts into a small number of latent dimensions (\textcite{gentzkow2019text}; \textcite{blei2003latent}). These approaches have been used successfully in a wide range of applications, including dictionary-based policy uncertainty indices (\textcite{baker2016measuring}), transcript-based political risk measures (\textcite{hassan2019firm}), and text-based measures of industry similarity (\textcite{hobergphillips2016textbased}). A common feature of these methods is that they summarize text through word frequencies or low-dimensional transformations of those frequencies, after which standard statistical tools are applied.

Recent advances in natural language processing instead rely on large language models (LLMs), which process text as sequences and form context-dependent representations of words and phrases. Modern LLMs are typically built on the transformer architecture (\textcite{vaswani2017attention}), which uses an attention mechanism to allow each word in a document to be interpreted in light of the surrounding text, including both nearby and more distant context. In addition, many contemporary models are instruction-tuned to follow natural-language task descriptions rather than merely predict the next word in a sequence (\textcite{ouyang2022training}). In our setting, we use such an instruction-following LLM in a prompt-based manner, treating it as a document annotator that receives a transcript and a task description and returns a structured assessment. This differs from traditional feature-based pipelines in that the mapping from text to output is generated directly by the model in response to the prompt, rather than through pre-specified word lists or estimated low-dimensional representations. We do not claim that this approach dominates alternative text-analysis techniques; rather, it reflects the view, supported by a growing literature on the use of LLMs in economics (\textcite{korinek2023generative}; \textcite{ludwig2025applied}), that context-sensitive, instruction-following models provide a natural additional tool for large-scale analysis of unstructured communications.
